\documentclass[10pt,a4paper]{amsart}
\usepackage[margin=19mm]{geometry}
\usepackage{amsmath,amssymb,mathtools}
\usepackage[scr=rsfs,cal=euler]{mathalfa}
\usepackage{xfrac}
\usepackage[unicode,hidelinks,bookmarks=false]{hyperref}
\newcommand{\ie}{i.e.\ }
\newcommand{\cf}{cf.\ }
\newcommand{\resp}{respectively\ }
\newcommand{\Subsection}{\subsection}
\providecommand{\nobreakdash}{\nobreak\hbox{-}}
\setlength{\emergencystretch}{2em}
\multlinegap0pt
\usepackage{float}
\usepackage{amssymb}
\usepackage[shortlabels]{enumitem}
\newtheorem{corollary}{Corollary}
\newtheorem{lemma}{Lemma}
\title{Characterizations of fractional operators via integral transforms}
\author{Daniel Cao Labora, Marc Jornet}
\date{Source: arXiv:2604.00714v1 (CC BY 4.0)}
\begin{document}
\maketitle

\noindent\emph{Source and licence.} Characterizations of fractional operators via integral transforms. Version arXiv:2604.00714v1 (CC BY 4.0), distributed by arXiv under the Creative Commons Attribution 4.0 International licence, http://creativecommons.org/licenses/by/4.0/, as recorded in the arXiv licence metadata for this exact version and shown on the abstract page https://arxiv.org/abs/2604.00714v1. Full licence notice: \texttt{licenses/arxiv-2604.00714-CC-BY-4.0.txt}.\par\smallskip

\noindent\textit{Editorial scope.} Counted unit: the article's lemma cfr (source0.tex lines 143--145). The article's complete original proof of it is delivered with it (source0.tex lines 146--156, one \texttt{proof} environment of 11 lines). No mathematics is written, repaired or re-derived: the statement and the proof are the article's own lines, in the article's own notation, and no retained line is substituted or re-flowed.  Retained uncounted as the source-local context the proof needs: lines 116--118; lines 135--136; lines 139--141.  Nothing was omitted from any retained span, so the package carries no omission record.  No retained line contains a citation, so the wrapper carries no bibliography and no bibliography entry.  No retained line contains a figure environment, an image-inclusion command or drawing code, so no image is delivered and the package declares no asset dependency.  Editorial formatting only, in addition: an AMS \texttt{amsart} preamble that restates the article's own declarations and macros used by the retained lines, namely the environments \texttt{corollary}, \texttt{lemma} on separate counters, together with the standard packages those lines need.  The retained environments print in this document as Corollary 1, Lemma 1 in order of appearance, whereas the article prints the same items with the numbering of its own full document.  The complete native source of the article as distributed in the arXiv e-print for this exact version is bundled unchanged as \texttt{sources/197606/source0.tex}.
\begin{equation} \label{Cauchy}
    h(x+y)=h(x)+h(y), \quad \forall x,y \in \Omega.
\end{equation}

\begin{corollary} \label{linear} Any function $h: \Omega=\prod_{j=1}^n (a_j,\infty) \subseteq (0,\infty)^n \to \mathbb{R}^n$ or $h: \Omega=\prod_{j=1}^n [a_j,\infty) \subseteq (0,\infty)^n \to \mathbb{R}^n$ that is continuous at one point and satisfies equation \eqref{Cauchy} has a unique extension to $\mathbb{R}^n$ that is continuous at one point and still satisfies equation \eqref{Cauchy}. In particular, the extension of $h$ is a linear function and, consequently, so is $h$.
\end{corollary}

\begin{equation} h(x+y)=h(x)+h(y),\quad \forall x,y\in\Omega\text{ such that }x+y\in\Omega. 
\label{cfd2}
\end{equation}

\begin{lemma} \label{cfr}
If $h:\Omega=(0,n)\rightarrow\mathbb{R}$ is a function that is continuous at one point and satisfies equation~\eqref{cfd2}, then $h$ is linear.
\end{lemma}

\begin{proof}
We extend $f$ to an additive function on $(0,\infty)$, so that Corollary~\ref{linear} can be applied.

Let $\eta\in [n,2n)$. We can write $\eta=\alpha+\beta$, where $0<\alpha<n$ and $0<\beta<n$. Then we define the extension $f(\eta):=f(\alpha)+f(\beta)$.

Let us see that this definition on $[n,2n)$ is well-posed. If $\eta$ were also $\tilde{\alpha}+\tilde{\beta}$, where $0<\tilde{\alpha}<n$ and $0<\tilde{\beta}<n$, then we consider without loss of generality that $\alpha<\tilde{\alpha}$ and $\tilde{\beta}<\beta$. By additivity on $(0,n)$, we know that $f(\tilde{\alpha})=f(\tilde{\alpha}-\alpha)+f(\alpha)$ and $f(\beta)=f(\beta-\tilde{\beta})+f(\tilde{\beta})$. Since $\beta-\tilde{\beta}=\tilde{\alpha}-\alpha$, this implies $f(\tilde{\alpha})+f(\tilde{\beta})=f(\alpha)+f(\beta)$, as desired.

Let us check that $f$ is additive on $[n,2n)$. Suppose that $\alpha\geq n$, $\beta<n$, and $\alpha+\beta<2n$. Since $0\leq \alpha-n<n-\beta$, we can pick a number $\alpha_2$ such that $\alpha-n<\alpha_2<n-\beta$, that is, $\alpha-\alpha_2<n$ and $\beta+\alpha_2<n$. Thus, by the definition of $f$, $f(\alpha+\beta)=f(\alpha-\alpha_2)+f(\beta+\alpha_2)$. On the one hand, $f(\beta+\alpha_2)=f(\beta)+f(\alpha_2)$, by additivity on $(0,n)$. On the other hand, again by the definition of $f$, $f(\alpha)=f(\alpha_2)+f(\alpha-\alpha_2)$. Therefore, we arrive at $f(\alpha+\beta)=f(\alpha)+f(\beta)$, as wanted.

Once we have the extension to $[n,2n)$, we can proceed analogously for $[2n,4n)$, $[4n,8n)$, etc.
\end{proof}

\end{document}
